% GENERATED FILE. Edit canonical lessons, then run npm run generate:books.
% canonical-source-hash: fnv1a64:ae6d3013ffebe99c
% canonical-lessons: FR-C421-annulation, FR-C421-remboursement, FR-C421-confirmer, FR-C421-decommander, FR-C421-surreservation

\chapter{Vol annulé}
\label{ch:fr-vol-annule}
\hlchaptermodality{421}

\begin{chapteropening}
\textbf{By the end of this chapter:} \emph{I can deal with a cancelled flight: the cancellation, the refund, confirming, calling off a booking and overbooking.}

Complete the last lesson of chapter 421: I can deal with a cancelled flight: the cancellation, the refund, confirming, calling off a booking and overbooking.
\end{chapteropening}

\section[l'annulation]{l'annulation — the cancellation}

\label{lesson:FR-C421-annulation}



\begin{quote}
Before the new one: say the French for the gap, then the French for to delay.
\end{quote}

\subsection*{You'll want to know: l'annulation}
\textbf{l'annulation} — \textquotedblleft{}the cancellation\textquotedblright{}.

\begin{quote}
Après l'annulation, la compagnie nous a donné une chambre.
\end{quote}

\begin{quote}
\emph{After the cancellation the airline gave us a room.}
\end{quote}

A feminine noun.

\begin{cousinweb}
\emph{annulation} is built on \emph{annuler} 'to cancel', a learned borrowing from Late Latin \emph{annullāre} 'to make nothing', from \emph{nullus} 'none' (Lewis \& Short).
\end{cousinweb}

\begin{grammarlens}[title={What you've built}]
One more word for telling, explaining and describing at B1.
\end{grammarlens}

\subsection*{Your turn}
\begin{itemize}
\raggedright
  \item \emph{Say it:} \emph{l'annulation}
  \item \emph{Say it:} \emph{l'annulation}, once more
  \item \emph{Say it:} a sentence of your own with \emph{l'annulation}
  \item \emph{Recall:} say the French for the gap, then the French for to delay, then say \emph{l'annulation} again
\end{itemize}

\begin{tcolorbox}[breakable,colback=teal!4,colframe=teal!35!black,title={Before you move on}]
What does l'annulation mean? (\textquotedblleft{}The cancellation\textquotedblright{}.) Say it once more.
\end{tcolorbox}
\section[le remboursement]{le remboursement — the refund}

\label{lesson:FR-C421-remboursement}



\begin{quote}
Before the new one: say the French for to delay, then the French for the cancellation.
\end{quote}

\subsection*{You'll want to know: le remboursement}
\textbf{le remboursement} — \textquotedblleft{}the refund\textquotedblright{}.

\begin{quote}
Elle attend le remboursement de son billet.
\end{quote}

\begin{quote}
\emph{She is waiting for the refund on her ticket.}
\end{quote}

Built on \emph{rembourser}.

\begin{cousinweb}
\emph{remboursement} is built on \emph{rembourser} 'to pay back', from \emph{bourse} 'purse', which comes from Old French \emph{borse} (EtymDB) and is usually traced to Late Latin \emph{bursa}, from Greek \emph{býrsa} 'hide'. EtymDB does not list the noun.
\end{cousinweb}

\begin{grammarlens}[title={What you've built}]
One more word for telling, explaining and describing at B1.
\end{grammarlens}

\subsection*{Your turn}
\begin{itemize}
\raggedright
  \item \emph{Say it:} \emph{le remboursement}
  \item \emph{Say it:} \emph{le remboursement}, once more
  \item \emph{Say it:} a sentence of your own with \emph{le remboursement}
  \item \emph{Recall:} say the French for to delay, then the French for the cancellation, then say \emph{le remboursement} again
\end{itemize}

\begin{tcolorbox}[breakable,colback=teal!4,colframe=teal!35!black,title={Before you move on}]
What does le remboursement mean? (\textquotedblleft{}The refund\textquotedblright{}.) Say it once more.
\end{tcolorbox}
\section[confirmer]{confirmer — to confirm}

\label{lesson:FR-C421-confirmer}



\begin{quote}
Before the new one: say the French for the cancellation, then the French for the refund.
\end{quote}

\subsection*{You'll want to know: confirmer}
\textbf{confirmer} — \textquotedblleft{}to confirm\textquotedblright{}.

\begin{quote}
L'agence a confirmé notre vol par courriel.
\end{quote}

\begin{quote}
\emph{The agency confirmed our flight by email.}
\end{quote}

\emph{La confirmation} — the confirmation.

\begin{cousinweb}
\emph{confirmer} is a learned borrowing from Latin \emph{cōnfirmāre} 'to make firm', from \emph{firmus} 'firm' (EtymDB).
\end{cousinweb}

\begin{grammarlens}[title={What you've built}]
One more word for telling, explaining and describing at B1.
\end{grammarlens}

\subsection*{Your turn}
\begin{itemize}
\raggedright
  \item \emph{Say it:} \emph{confirmer}
  \item \emph{Say it:} \emph{confirmer}, once more
  \item \emph{Say it:} a sentence of your own with \emph{confirmer}
  \item \emph{Recall:} say the French for the cancellation, then the French for the refund, then say \emph{confirmer} again
\end{itemize}

\begin{tcolorbox}[breakable,colback=teal!4,colframe=teal!35!black,title={Before you move on}]
What does confirmer mean? (\textquotedblleft{}To confirm\textquotedblright{}.) Say it once more.
\end{tcolorbox}
\section[décommander]{décommander — to cancel (a booking, an order)}

\label{lesson:FR-C421-decommander}



\begin{quote}
Before the new one: say the French for the refund, then the French for to confirm.
\end{quote}

\subsection*{You'll want to know: décommander}
\textbf{décommander} — \textquotedblleft{}to cancel (a booking, an order)\textquotedblright{}.

\begin{quote}
Après l'annulation du vol, nous avons décommandé le taxi.
\end{quote}

\begin{quote}
\emph{After the flight was cancelled, we cancelled the taxi.}
\end{quote}

\emph{Se décommander} — to cry off, to cancel at the last minute.

\begin{cousinweb}
\emph{décommander} is built from \emph{dé-} and \emph{commander} 'to order', which goes back through Vulgar Latin to Latin \emph{commendāre} 'to entrust', built on \emph{mandāre} (EtymDB lists \emph{commander}).
\end{cousinweb}

\begin{grammarlens}[title={What you've built}]
One more word for telling, explaining and describing at B1.
\end{grammarlens}

\subsection*{Your turn}
\begin{itemize}
\raggedright
  \item \emph{Say it:} \emph{décommander}
  \item \emph{Say it:} \emph{décommander}, once more
  \item \emph{Say it:} a sentence of your own with \emph{décommander}
  \item \emph{Recall:} say the French for the refund, then the French for to confirm, then say \emph{décommander} again
\end{itemize}

\begin{tcolorbox}[breakable,colback=teal!4,colframe=teal!35!black,title={Before you move on}]
What does décommander mean? (\textquotedblleft{}To cancel (a booking, an order)\textquotedblright{}.) Say it once more.
\end{tcolorbox}
\section[la surréservation]{la surréservation — the overbooking}

\label{lesson:FR-C421-surreservation}



\begin{quote}
Before the new one: say the French for to confirm, then the French for to cancel.
\end{quote}

\subsection*{You'll want to know: la surréservation}
\textbf{la surréservation} — \textquotedblleft{}the overbooking\textquotedblright{}.

\begin{quote}
À cause de la surréservation, deux passagers n'ont pas pu partir.
\end{quote}

\begin{quote}
\emph{Because of the overbooking, two passengers could not leave.}
\end{quote}

Many people say \emph{le surbooking}.

\begin{cousinweb}
\emph{surréservation} is built from \emph{sur-} 'over' and \emph{réservation}, from \emph{réserver} 'to keep aside', a learned borrowing from Latin \emph{reservāre} (Lewis \& Short). EtymDB does not list it.
\end{cousinweb}

\begin{grammarlens}[title={What you've built}]
One more word for telling, explaining and describing at B1.
\end{grammarlens}

\subsection*{Your turn}
\begin{itemize}
\raggedright
  \item \emph{Say it:} \emph{la surréservation}
  \item \emph{Say it:} \emph{la surréservation}, once more
  \item \emph{Say it:} a sentence of your own with \emph{la surréservation}
  \item \emph{Recall:} say the French for to confirm, then the French for to cancel, then say \emph{la surréservation} again
\end{itemize}

\begin{tcolorbox}[breakable,colback=teal!4,colframe=teal!35!black,title={Before you move on}]
What does la surréservation mean? (\textquotedblleft{}The overbooking\textquotedblright{}.) Say it once more.
\end{tcolorbox}
